\documentclass[../main.tex]{subfiles}
\begin{document}
\chapter{Series of Functions}
\section{Definitions}
In the previous chapter, we extended our theory of sequences of numbers to sequences of functions.
We would also like to extend the idea of series of numbers to series of functions.
\begin{definition}[$n$-th Partial Sum]
  Let $g_n: E \to \R$ be  a sequence of functions.
  We define its \textit{$n$-th partial sum} to be the sequence of functions $f_n: E \to \R$ given by:
  \[
    f_n(x) = \sum_{j = 1}^{n} g_j(x)
  \]
\end{definition}
\begin{definition}[Convergence of a series of functions at a point]
  We say that the series of functions $\sum_{j = 1}^{\infty} g_j$ converges at a point $x \in E$ if the sequence $(f_n(x))$ converges.
\end{definition}
\begin{definition}[Uniform convergece for series of functions]
  We say that the series of functions $\sum_{j = 1}^{\infty} g_j$ \textit{converges uniformly} on $E$ if $(f_n)$ converges uniformly on $E$.
\end{definition}
\section{Absolute Convergence}
We can also extend the concept of absolute convergence.
\begin{definition}[Absolute convergence for series of functions at a point]
  We say that the series of functions $\sum_{j = 1}^{\infty} g_j$ \textit{converges absolutely} at a point $x \in E$ if the series $\sum_{j = 1}^{\infty} |g_j(x)|$ converges.
\end{definition}
\begin{definition}[Absolute uniform convergence for series of functions]
  We say that the series of functions $\sum_{j = 1}^{\infty} g_j$ \textit{converges absolutely uniformly} on $E$ if the series of functions $\sum_{j = 1}^{\infty} |g_j|$ converges uniformly.
\end{definition}
\begin{remark}[Recap]
  Recall from Analysis I that for series of numbers in $\C$, absolute convergence of the series implies convergence.
\end{remark}
\begin{proposition}
  If $g_j: E \to \R$ is a sequence of functions and $\sum_{j = 1}^{\infty} g_j$ converges uniformly absolutely on $E$, then $\sum_{j = 1}^{\infty} g_j$ converges uniformly on $E$.
\end{proposition}
\begin{proof}
  Define the partial sums $f_n(x) = \sum_{j = 1}^{n} g_j(x)$ and $h_n(x) = \sum_{j = 1}^{n} |g_j(x)|$.
  We do not know what the limit function of $f_n$ should be and so we want to appeal to \cref{cauchyCritereon}.
  Considering $n > m$, we have:
  \begin{align*}
    |f_n(x) - f_m(x)| &= \abs{\sum_{j = m + 1}^{n} g_j(x)} \\
                      &\leq \sum_{j = m + 1}^{n} |g_j(x)| \\
                      &= h_n(x) - h_m(x) \\
                      &= |h_n(x) - h_m(x)|
  \end{align*}
  Therefore:
  \[
    \sup_{x \in E} |f_n(x) - f_m(x)| \leq \sup_{x \in E} |h_n(x) - h_m(x)|
  \]
  Since $\sum_{j = 1}^{\infty} g_j$ converges absolutely uniformly, $h_n$ converges uniformly and so by \cref{cauchyCritereon} it is uniformly Cauchy.
  Thus as $m \to \infty$ $\sup_{x \in E} |f_n(x) - f_m(x)| \to 0$ and so $(f_n)$ converges uniformly by \cref{cauchyCritereon}.
\end{proof}
\begin{remark}
  Uniform convergence of $\sum_{j = 1}^{\infty} g_j$ and absolute convergence of $\sum_{j = 1}^{\infty} g_j(x)$ at every $x \in E$ is \textbf{not} enough to conclude that a series is absolutely uniformly convergent.

  For example, consider the series of functions defined on $E = (0, 1)$:
  \[
    \sum_{j = 1}^{\infty} \frac{(-1)^j}{j} x^j
  \]
  Using the ratio test for power series from IA Analysis I, $\sum_{j = 1}^{\infty} \frac{x^j}{j}$ has radius of convergence $R = 1$.
  Thus $\sum_{j = 1}^{\infty} \frac{(-1)^j}{j} x^j$ converges absolutely as a series of numbers at every $x \in (0, 1)$.

  We will see on Example Sheet 1 Q7 that the series converges uniformly on $(0, 1)$.
  You will see that we are able to do this because the series $\sum_{j = 1}^{\infty} \frac{(-1)^j}{j}$ converges (by the alternating series test) which is the behaviour at the endpoint of $(0, 1)$.

  However, when we consider the absolute version, we now have that $\sum_{j = 1}^{\infty} \frac{1}{j}$ at the endpoint and so we cannot use the strategy above.
  In fact, this does not converge absolutely uniformly on $(0, 1)$.
  We can demonstrate this by showing that it violates the uniform Cauchy criterion.
  For $n > m$ we have:
  \begin{align*}
    \sum_{j = 1}^{n} \abs{\frac{(-1)^j}{j} x^j} - \sum_{j = 1}^{m} \abs{\frac{(-1)^j}{j} x^j} &= \sum_{j = n + 1}^{n} \frac{x^j}{j} \\
                                                                                              &\geq \sum_{j = m + 1}^{n}\frac{x^n}{j}
  \end{align*}
  Taking the supremum of both sides over $(0, 1)$:
  \[
    \sup_{x \in (0, 1)} \left[\sum_{j = 1}^{n} \abs{\frac{(-1)^j}{j} x^j} - \sum_{j = 1}^{m} \abs{\frac{(-1)^j}{j} x^j}\right] \geq \sum_{j = m + 1}^{n} \frac{1}{j}
  \]
  But for any fixed $m$, as $n \to \infty$, $\sum_{j = m + 1}^{n} \frac{x^{n}}{j} \to \infty$ and so $\sum_{j = 1}^{\infty} \abs{\frac{(-1)^j}{j} x^j}$ does not converge uniformly on $(0, 1)$.
\end{remark}
We have the following simple but useful criterion for absolute uniform convergence.
\begin{theorem}[Weierstrass $M$-test]
  Let $g_j: E \to \R$ be a series of functions.
  If there exists $M_j \geq 0$ such that:
  \[
    \sup_{x \in E} |g_j(x)| \leq M_j\ \forall j \in \N
  \]
  and $\sum_{j = 1}^{\infty} M_j$ converges, then $\sum_{j = 1}^{\infty} g_j$ converges absolutely uniformly on $E$.
\end{theorem}
\begin{proof}
  Define the partial sums $h_n(x) = \sum_{j = 1}^{n} |g_j(x)|$ and $s_n = \sum_{j = 1}^{n} M_j$.

  We again wish to appeal to the Cauchy critereon and so considering $n > m$:
  \begin{align*}
    h_n(x) - h_m(x) &= \sum_{j = m + 1}^{n} |g_j(x)| \\
                    &\leq \sum_{j = m + 1}^{n} M_j \\
                    &= s_n - s_m
  \end{align*}
  Thus for all $n, m$:
  \[
    \sup_{x \in E} \abs{h_n(x) - h_m(x)} \leq |s_n - s_m|
  \]
  Since $\sum_{j = 1}^{\infty} M_j$ converges, its partial sums $(s_n)$ are Cauchy and so $h_n$ is uniformly Cauchy on $E$.
  Thus by \cref{cauchyCritereon}, $h_n$ converges uniformly on $E$, i.e. $\sum_{j }^{\infty} g_j$ converges uniformly absolutely on $E$.
\end{proof}
\section{Power Series}
\end{document}
